\section{Critical roofs, actual reversing orbits, and their exact density coefficient}
\label{sec:cr-critical}
We now identify a dynamical invariant in the critical-roof density,
without averaging the roof or replacing the physical trajectories.
Let $x$ be a regular $m$-flight trajectory with endpoints in the
original section cap. Write $q_x=(\alpha_0,\alpha_m)$ and
$t_x=F(q_x)$. All statements about derivatives of the actual
collision map below concern regular trajectories, not a formal
continuation through a tangent collision.

\begin{theorem}[Critical half-paths are physical reversing periodic orbits]
\label{thm:cr-periodic}
The endpoint roof is critical precisely when $p_0=p_m=0$.
Such a trajectory extends by the actual reflection law to a regular
$2m$-collision periodic orbit, whose period need not be primitive.
The extra $m$ flights traverse the original flights in reverse,
with total roof $2t_x$ and total lifted displacement zero.
Let $\mathcal I(\alpha,p)=(\alpha,-p)$ be collision time reversal,
and put
\[
 A^\mathcal I_{m,R}(x)
       =\sum_{j=0}^{m-1}\one_{Y_R^*}(\mathcal I T_R^j x).
\]
If the half-path belongs to the original exact event, then
\begin{equation}\label{eq:cr-occupation}
 A_{2m,R}(x)=A_{m,R}(x)+A^\mathcal I_{m,R}(x).
\end{equation}
One cannot replace the right side by $2n$ without additionally
verifying the time-reversed section tests.
\end{theorem}
\begin{proof}
The first-variation identity already retained in
\eqref{eq:v109-density} is $\nabla_qF=R(-p_0,p_m)$, which proves
the criticality assertion. At a normal terminal collision the
reflection sends the incoming velocity to its negative. The ray
therefore retraces the preceding actual flight. At each intervening
regular collision the reversed reflection retraces the next
preceding flight. At the normal initial collision it returns to the
original outgoing state. Thus every reversed flight is an actual
free flight, with no extra or hidden collision, and the orbit closes
after $2m$ collisions. Its reversed displacement is the negative of
the forward displacement. Its length is twice the original length.
No assertion about a primitive period is used.

For $1\le j<m$ the reversed state at time $2m-j$ is
$\mathcal I T_R^j x$, while the states at times $0$ and $m$ are
fixed by $\mathcal I$. In the original event both endpoint section
tests equal one. Summing the $m$ forward tests and then the terminal
test and the $m-1$ reversed interior tests gives
\eqref{eq:cr-occupation}, since the initial test in
$A^\mathcal I_m$ also equals one. The original section is not
assumed invariant under $\mathcal I$.
\end{proof}

\begin{theorem}[Exact stability determinant at a critical roof]
\label{thm:cr-monodromy}
At a regular critical half-path write
\[
 G_x=\begin{pmatrix}a&e\\e&d\end{pmatrix},\qquad
 \Delta_x=ad-e^2>0,\qquad
 \mathcal M_x=D T_R^{2m}(x).
\]
Then $e\ne0$, and the actual return derivative satisfies
\begin{equation}\label{eq:cr-monodromy}
 \det\mathcal M_x=1,\qquad
 \operatorname{tr}\mathcal M_x-2
   =|\det(I_2-\mathcal M_x)|
   =\frac{4\Delta_x}{e^2}>0.
\end{equation}
Consequently the critical density coefficient
\begin{equation}\label{eq:cr-jump}
 \mathfrak J(x):=\frac{\omega(x)}{2c_*\sqrt{\det G_x}}
   =\frac{1}{Rc_*\sqrt{|\det(I_2-D T_R^{2m}(x))|}}
\end{equation}
is an invariant of the actual reversing periodic orbit with its
marked initial state and specified half-length $m$. The coefficient
is not a probability conditioned on that periodic orbit.
\end{theorem}
\begin{proof}
The signed optical corner formula \eqref{eq:v110-signed-corner}
makes $e$ nonzero on every regular finite trajectory. Differentiate
$p_0=-F_{\alpha_0}/R$ and $p_m=F_{\alpha_m}/R$. Solving the first
identity for $d\alpha_m$ gives, in the original $(\alpha,p)$
coordinates,
\begin{equation}\label{eq:cr-half-derivative}
 B_x:=D T_R^m(x)
  =\begin{pmatrix}
       -a/e&-R/e\\-\Delta_x/(Re)&-d/e
    \end{pmatrix},\qquad \det B_x=1.
\end{equation}
This keeps the sign of $e$; replacing it prematurely by $|e|$ in
the matrix would change a derivative, although not its final trace.
At both normal endpoints the derivative of time reversal is
$J=\operatorname{diag}(1,-1)$. The reversibility identity
$\mathcal I T_R^m\mathcal I=T_R^{-m}$ therefore yields
$D T_R^m(T_R^m x)=J B_x^{-1}J$ and
\[
 \mathcal M_x=J B_x^{-1}J B_x.
\]
For an arbitrary determinant-one matrix
$B=\left(\begin{smallmatrix}b_{11}&b_{12}\\b_{21}&b_{22}\end{smallmatrix}\right)$,
$J B^{-1}J=\left(\begin{smallmatrix}b_{22}&b_{12}\\b_{21}&b_{11}\end{smallmatrix}\right)$,
so
\[
 \operatorname{tr}(J B^{-1}J B)-2=4b_{12}b_{21}.
\]
Substitution of \eqref{eq:cr-half-derivative} gives
$4\Delta_x/e^2$. For a two-dimensional determinant-one matrix,
$\det(I_2-\mathcal M_x)=2-\operatorname{tr}\mathcal M_x$.
This proves \eqref{eq:cr-monodromy} and also gives positive real
hyperbolic multipliers. Finally $\omega=|e|/R$, so elementary
simplification proves \eqref{eq:cr-jump}. The determinant of
$I_2-\mathcal M_x$ is invariant under a smooth change of collision
coordinates. No periodic-orbit trace formula or distributional
limit theorem has been invoked.
\end{proof}

\begin{theorem}[Critical-roof coarea with arbitrary exact-event cuts]
\label{thm:cr-critical-density}
Fix a regular critical half-path and a sufficiently small endpoint
neighbourhood $V$ on its designated itinerary. Let $E\subset V$ be
an arbitrary Borel selector, including the actual physical domain,
initial and terminal section tests, exact occupation and displacement
labels, and any clearance tests. Put
\[
 y=G_x^{1/2}(q-q_x),\quad f(y)=F(q)-t_x,\quad
 w(y)=\frac{\omega(q)}{4\pi c_*\sqrt{\det G_x}}.
\]
Choose $r_0>0$ so that the closed $y$-ball is contained in the
coordinate neighbourhood, and constants $K,L_w\ge0$ such that
\[
 \|D^2 f(y)-I_2\|\le K|y|,\quad
 |w(y)-w(0)|\le L_w|y|,\quad Kr_0\le1.
\]
For $0<s<r_0^2/3$ there is, for every angle $\theta$, a unique
$r_s(\theta)\in(0,r_0)$ with $f(r_s e(\theta))=s$. Define
\[
 a_E(s)=\frac{1}{2\pi}\int_0^{2\pi}
     \one_E\bigl(q_x+G_x^{-1/2}r_s(\theta)e(\theta)\bigr)\,d\theta.
\]
The original positive collision measure restricted to $E$ and this
ball has a density $D_E$ near $t_x$, no atom at $t_x$, and
\begin{equation}\label{eq:cr-sector}
 \left|D_E(t_x+s)-\mathfrak J(x)a_E(s)\right|
 \le 2\pi\sqrt3\{2L_w+Kw(0)\}\sqrt{s}.
\end{equation}
There is no contribution below $t_x$. The constant is independent
of the Borel selector $E$, but is a constant of this local chart,
not a claimed uniform-in-$m$ constant.

If $E$ contains a neighbourhood of $q_x$, then $a_E(s)=1$ for
small $s$ and the one-sided density limit is exactly
$\mathfrak J(x)$. If the selected physical side is described by
finitely many $C^1$ inequalities vanishing at $q_x$ with nonzero
linear parts, then its limiting coefficient is
$\mathfrak J(x)$ times the angular fraction of the corresponding
linear sector. Without such a hypothesis the exact $a_E(s)$ is
retained; no limit for it is asserted.
\end{theorem}
\begin{proof}
The retained reduced action is smooth on a regular finite itinerary;
the inverse function theorem and the nonzero optical corner give a
smooth positive $\omega$ there. Thus the displayed constants exist
after decreasing $r_0$. At a critical point $f(0)=0$,
$\nabla f(0)=0$ and $D^2f(0)=I_2$. Integrating the Hessian bound
along a radial segment gives
\[
 |\partial_r f(re)-r|\le Kr^2/2,\qquad
 |f(re)-r^2/2|\le Kr^3/6.
\]
For $r\le r_0$ these imply
$r^2/3\le f(re)\le2r^2/3$, $\partial_r f\ge r/2$, and hence
the existence and uniqueness of $r_s$, with $r_s\le\sqrt{3s}$.
They also give
\[
 \left|\frac{r_s}{\partial_r f(r_s e)}-1\right|\le Kr_s,
 \qquad 0<\frac{r_s}{\partial_r f(r_s e)}\le2.
\]

In endpoint coordinates the original collision density is
$\omega(q)/(4\pi c_*)$; the linear substitution to $y$ makes it
$w(y)\,dy$. Polar coordinates and the monotone radial substitution
$t=t_x+f(re(\theta))$, tested against nonnegative roof functions,
give the exact density
\[
 D_E(t_x+s)=\int_0^{2\pi}\one_E(q_s(\theta))
       w(r_s e(\theta))
       \frac{r_s}{\partial_r f(r_s e(\theta))}\,d\theta.
\]
Here $q_s(\theta)=q_x+G_x^{-1/2}r_s e(\theta)$, and any source
outside the selected ball is not part of $D_E$. The point $r=0$
has zero two-dimensional measure, so it produces no atom. The
strict radial positivity also excludes mass below $t_x$ inside the
ball. For each angle the difference between the weight in this
integral and $w(0)$ is at most
$(2L_w+Kw(0))r_s$. Integrating gives
\eqref{eq:cr-sector}, since $2\pi w(0)=\mathfrak J(x)$.
This estimate does not differentiate an indicator or require its
boundary to have zero roofwise length.

When $E$ contains a neighbourhood, every sufficiently small radial
point belongs to it. For the stated finite collection of $C^1$
cuts, Taylor expansion of each defining function shows convergence
of its sign to its nonzero linear part at every angle except the
finitely many zero directions. Dominated convergence gives the
sector fraction. It is the fraction after the transformation by
$G_x^{-1/2}$, not an unweighted fraction in the original endpoint
angles. For a general Borel selector no such argument is used.
\end{proof}

\begin{corollary}[One-sided unvisited boundaries without a fictitious orbit]
\label{cor:cr-one-sided}
Suppose a designated finite broken path reaches an unvisited-disk
tangency or a section cut, all its \emph{visited} incidences remain
positive, and its endpoint action is critical. The sector formula
\eqref{eq:cr-sector}, with the original physical-side selector,
continues to hold in the one-sided coordinate neighbourhood.
At a singular state the second expression in \eqref{eq:cr-jump}
is \emph{not} asserted to be a derivative of a defined physical
return map; the first expression is the one-sided action coefficient.
\end{corollary}
\begin{proof}
The stationary equations for the visited interior collision points
have positive-definite interior Hessian. At the stated limiting path
they therefore have a smooth local solution as functions of the two
endpoints. The signed optical corner is nonzero because the
visited incidences and flight lengths are positive. This supplies
the smooth action and endpoint-coordinate Jacobian on a full
auxiliary neighbourhood. The portion representing actual trajectories
is retained as a Borel subset $E$; portions intersecting a forbidden
disk are never assigned to the physical event. On the admitted side
these functions agree with the actual action and original measure.
The proof of Theorem~\ref{thm:cr-critical-density} used arbitrary
Borel restrictions, so it applies to this subset without continuing
the physical dynamics across its boundary. Simultaneous cuts are
included in the same selector. If visited grazing also occurs, this
particular smooth extension is not invoked. Its treatment remains
the inherited visited-grazing trace argument. Nothing in this
corollary equates an area-null boundary point with a zero positive
pairing-endpoint trace.
\end{proof}

\begin{proposition}[A complete finite-count critical coefficient bound]
\label{prop:cr-critical-sum}
Let the sum run over all distinct regular critical half-paths in
the original exact event, at fixed $m,R,n,k$. Then
\begin{equation}\label{eq:cr-critical-one}
 \mathfrak J(x)\le
     \frac{47}{12c_*\lambda_\circ}\varrho^{m-1}.
\end{equation}
Under the retained, explicitly stated geometric-pressure input
\ref{input:v112-geometric-pressure}, the full positive sum satisfies
\begin{equation}\label{eq:cr-critical-sum}
 \sum_x\mathfrak J(x)
          \le\frac{C_0Q_*}{2c_*\lambda_\circ}.
\end{equation}
This is unscaled. It does not supply an $m^{-2}$ bound for a
clearance-selected critical source or a uniform bound for the local
error constants in \eqref{eq:cr-sector}.
\end{proposition}
\begin{proof}
Use $\sqrt{\det G_x}\ge\lambda_\circ$ in
\eqref{eq:cr-jump} and the optical bound for $\omega$ in
\eqref{eq:v109-optical-bounds} to obtain
\eqref{eq:cr-critical-one}. The retained global equal-gradient
uniqueness gives at most one regular critical half-path per lifted
physical word. Such a path extends periodically without singular
collisions, so it belongs to the two-sided regular set $M_R'$.
The selected-word sum in \eqref{eq:v111-selected-word-sum}, with
no incidence marks, therefore bounds $\sum_x\omega(x)$ by
$C_0Q_*$. Divide by $2c_*\lambda_\circ$. Missing words contribute
zero. No multiplication of a per-orbit exponential estimate by the
number of words is made.
\end{proof}

\subsection{What the critical calculation does and does not close}
Formula \eqref{eq:cr-jump} identifies an actual return-stability
weight in the original roof density; it is not an ansatz for a
Markov model and not a periodic-orbit replacement of the invariant
measure. Formula \eqref{eq:cr-sector} also keeps positive source on
section-cut and one-sided unvisited boundaries rather than assigning
it a zero coefficient by area-nullness. The full paired source,
its nongrazing comparison pool, and \eqref{eq:v112-remaining} are
unchanged. A critical \emph{component} coefficient need not equal
a discontinuity of the entire density, since other physical
components may be present at the same roof.

For a collision-uniform payment it is still necessary to control
the actual weighted sum on fixed outer neighbourhoods, its
clearance and exact-label localization, and the chart remainders.
Choosing a separate collision-dependent neighbourhood merely to
separate critical roof values would not prove that fixed-window
statement. Neither \eqref{eq:cr-critical-one} nor the unscaled
sum \eqref{eq:cr-critical-sum} is used in that invalid manner.
The new identities instead give the exact geometric and stability
coefficients that a positive capacity or transfer-operator estimate
must control.
